\Chapter{1}

\Verse{1}
{ \hX{} }
{ \eX{} }
{
    I typed too much so now the stuff below needs an explanation.

    Let's make the explanation be that it's too long because it's one of those ``epistle'' things.

    Like the Paul shit, where he goes on and on and on forever and it never ends.

    And then eventually someone puts his emails to various people in the bible for some reason, and barely anyone reads them, cuz (again) they're way too long.
}

\Verse{2}
{ \hX{} }
{ \eX{} }
{
    On that note, we now begin a very long letter by some a\eR{po}s\eR{tle or }s\eR{aint or }hole\eR{y individual} named J, written to an unnamed church in Idaho.

    \emph{(ahem)}
}

\Verse{3}
{ \hX{} }
{ \eX{} }
{
    1 Jeffesians.

    1 Jeffesians 1.

    An epistle of J, apostle of J, from the International House of παν Cakes,\textsuperscript{0} by the will of God,

    To Jeff of the church of Jeffesians in Idaho.

    \emph{(Continued from previous email.)}
}

\Verse{4}
{ \hX{} }
{ \eX{} }
{
    Dear Jeffesians,

    As you can tell from my preamble it's gonna be quite a long letter.
}

\Verse{5}
{ \hX{} }
{ \eX{} }
{
    But yeah if we didn't even know \emph{that} basic thing about the grammar of CBH\textsuperscript{1} until our lifetimes, the amount of low hanging fruit\textsuperscript{2} and easily solvable problems in the field is likely to be huge given our current knowledge from source analysis and archaeology.
}

\Verse{6}
{ \hX{} }
{ \eX{} }
{
    Which is wild because most people think there's not much ``new'' to be learned about a 3000 year old text.

    Totally wrong.

    Sure people have been reading the damn thing for 3000 years.

    But it's only in our lifetime that we've known enough to start understanding what the authors were up to.
}

\Verse{7}
{ \hX{} }
{ \eX{} }
{
    I've been in a lot of interesting fields, but it's hard to find interesting new solvable problems in (say) theoretical physics, and other fields where new solvable problems are more abundant like evolutionary psychology, you mostly have to undertake that new research inside of institutions that, even when I left like 10 years back, were already so far beyond the era when they had been\textsuperscript{3} were functional and truth-seeking that even with the best mentors in the field the majority of your time is wasted on institutional friction (e.g., department shit, experiment review board shit, academic journal shit, and all the other stuff the old Ivory Tower has since become.)
}

\Verse{8}
{ \hX{} }
{ \eX{} }
{
    End of rant.

    The point of the above is that bible research is different.

    \begin{enumerate}

    \item It's easier to make new contributions to unsolved problems in biblical text criticism and source analysis than it is in (a) mathematics, (b) theoretical physics, and (c) experimental psychology.

    \item There's no institutional bullshit, because the bible's open source (thanks Martin Luther you wild bastard)

    \item Like half the planet belongs to some religion that takes the old testament as part of its text, so half the planet sort of feels like they \emph{should} read and understand the book, but

    \item All existing translations and commentaries are dogshit. Even the best ones. Even the ones by my personal heroes like Richard Elliott Friedman and Baruch Halpern. Those dudes made massive contributions to the field that the current book (and forthcoming books) would never have be possible without. I could never have done what those guys did.

    \item But they're both academics, and you can't write a book like this while you're part of academia.

    \end{enumerate}
}

\Verse{9}
{ \hX{} }
{ \eX{} }
{
    The bible breaks too many rules.

    \begin{itemize}

    \item It doesn't cite sources.

    \item It has no copyright page.

    \item It has no named author.

    \item It has too much. (Duplicates of tons of stories that any modern editor would cut.)

    \item It has too little. (It doesn't explain itself or help the reader out in any way)

    \item No single person gets the credit for it (even a fictitious person.)

    \item No single person gets the royalties for it.

    \end{itemize}
}

\Verse{10}
{ \hX{} }
{ \eX{} }
{
    And most importantly:

    \begin{itemize}

    \item Anyone can steal it and make their own version, cuz it's open source.

    \end{itemize}
}

\Verse{11}
{ \hX{} }
{ \eX{} }
{
    But even though it's open source, the existing versions are entirely uncreative in how they present the damn text.

    Basically all existing translations and commentaries either came out of:

    \begin{itemize}

    \item Religion, or

    \item Academia

    \end{itemize}
}

\Verse{12}
{ \hX{} }
{ \eX{} }
{
    Which are two places where you \emph{absolutely} cannot be as weird and silly as the text itself.

    Which means you can't show a modern reader what the text \emph{is,} even in a totally non-ideological way.

    Which means \emph{no previous version has even tried to represent what's actually in the damn book.}
}

\Verse{13}
{ \hX{} }
{ \eX{} }
{
    They either write books \emph{about} the bible (in the genre of ``telling you about it,'' but not being an actual version of IT.)

    Or they write translations and commentaries, either from inside a church affiliated institution, or an academic one.
}

\Verse{14}
{ \hX{} }
{ \eX{} }
{
    And you're not allowed to say ``dick'' in either of those places.

    And you're certainly not allowed to invent a new genre and then not explain what you're doing, but let the genre explain itself.

    BUT THAT'S WHAT THE DAMN BIBLE DOES!
}

\Verse{15}
{ \hX{} }
{ \eX{} }
{
    In summary:

    \emph{It's illegal to write a bible in the modern world.}
}

\Verse{16}
{ \hX{} }
{ \eX{} }
{
    If you try to BE the genre, not just explain it, then:

    \begin{itemize}

    \item You must not cite sources.

    \item You can't put your name on it.

    \item You can't have an author's Preface or Introduction that explains what the book is.

    \item You can't put urls or contact information in the book, unless it's the information of the publisher.

    \item If you do that, then the publisher can't be identified as the author.

    \item All of this is simple kindergarten ``what is the bible'', and basically 0\% of it is stuff that modern publishers are ok with, whether inside academia or in the popular press.

    \end{itemize}
}

\Verse{17}
{ \hX{} }
{ \eX{} }
{
    Anyways, in summary:

    It's illegal to write a bible.

    It is therefore necessary that we do it.
}

\Verse{18}
{ \hX{} }
{ \eX{} }
{
    In the name of the Genre, the Authors, and the Text.

    \noindent z\\
    a \chinese{咱們}\textsuperscript{4} men\\
    n
}

\Verse{19}
{ \hX{} }
{ \eX{} }
{
    [0: παν. Greek \emph{pan}, meaning ``all.'' The location of the biblical International House Of All Cakes is unknown.]

    [1: Classical Biblical Hebrew.]

    [2: Pun intended.]

    [3: An example of the past perfect tense in Classical Biblical English.]

    [4: The term zanmen (\chinese{咱們}) means ``We''. It is an inclusive we (i.e., ``you and me'', as in ``we have fun together'') rather than an exclusive we (i.e., me and someone else'', as in ``we think you have a problem.'']
}

\Verse{20}
{ \hX{} }
{ \eX{} }
{
    This ends J's epistle to the church of the Jeffesians in Idaho that eventually moved south, to the land of Edom, where the bible is read.
}
